\section{Setting and Notation}
\label{sec:setting}

\paragraph{Fields and cell ring.} Fix
$P(X) = X^{192} + X^7 + X^2 + X + 1 \in \F_2[X]$ — the NIST FIPS 186-2
pentanomial — and write
\[
  \GFt \;:=\; \F_2[X] \big/ \bigl(P(X)\bigr)
\]
for the binary field of size $2^{192}$. Fix a coefficient bound
$\Dval = 32$ and define the \emph{cell ring}
\[
  \cellring \;:=\; \F_2[X]\big/(X^{\Dval}),
\]
a free $\F_2$-module of rank $\Dval$. We treat $\cellring$ as a
subset of $\polyringat{\Dval} \subset \polyring$ throughout; the
quotient $/(X^\Dval)$ is structural (it bounds the cell size) and is
\emph{not} the source of any reduction in this document.

\paragraph{The two roles of $X$.} The formal variable $X$ plays two
distinct roles in the construction:
\begin{itemize}[leftmargin=2em]
  \item As the \emph{coefficient variable} of cells in $\cellring$
    and of intermediate polynomials in $\polyring$. Operations on
    cells (multiplication, evaluation) live here.
  \item As the \emph{quotient generator} of $\GFt$. An element
    $g \in \GFt$ has a unique degree-$<192$ representative
    $g_{\mathrm{rep}}(X) \in \polyringat{192}$, the canonical lift.
\end{itemize}
The two roles are notationally the same letter $X$; they are
distinguished structurally by the surrounding ring. When confusion is
possible we write $g_{\mathrm{rep}}$ or $\polyringat{192}$ explicitly.

\paragraph{Trace and witness.} Fix a row count $n = 2^{\nu}$ and a
column count $K$. The witness is $K$ trace columns
$w_1, \ldots, w_K \in \cellring^n$, each a length-$n$ vector of cells.
For analysis we will arrange a single column $w$ as an
$n_{\mathrm{rows}} \times n_{\mathrm{cols}}$ matrix
$M_w$ with $n = n_{\mathrm{rows}} \cdot n_{\mathrm{cols}}$, both
factors powers of two; the row index splits as
$b = (b_0, b_1) \in \{0,1\}^{\nu_0} \times \{0,1\}^{\nu_1}$ with
$\nu = \nu_0 + \nu_1$.

\paragraph{Evaluation projection $\psia$.} Fix a transcript-fresh
$\alpha \in \GFt$, drawn after the witness is committed. Define the
$\F_2$-algebra homomorphism
\[
  \psia : \polyring \;\longrightarrow\; \GFt,
  \qquad
  \psia\!\big(\textstyle\sum_j c_j X^j\big) \;=\; \sum_j c_j \cdot \alpha^j \pmod{P}.
\]
Restricted to $\cellring \subset \polyringat{\Dval}$, $\psia$ is the
PIOP's Step~3 projection — it replaces each cell by a single
$\GFt$-element. We use the same symbol $\psia$ for the cell-wise
application $\psia : \cellring^n \to \GFt^{\,n}$.

\paragraph{Multilinear extension.} For $v \in \cellring^n$, the
\emph{multilinear extension over $\GFt$} of $\psia(v)$ is the
unique multilinear polynomial $\mle{\psia(v)} \in \GFt[Y_1,\ldots,Y_\nu]$
with $\mle{\psia(v)}(b) = \psia(v_b)$ for every $b \in \{0,1\}^\nu$.
Equivalently,
\[
  \mle{\psia(v)}(r) \;=\; \sum_{b \in \{0,1\}^\nu} \eq_b(r) \cdot \psia(v_b),
  \qquad
  \eq_b(r) := \prod_{k=1}^\nu \bigl(b_k r_k + (1 - b_k)(1 - r_k)\bigr) \in \GFt.
\]

\paragraph{The Step-7 input.} The protocol's pre-Step-7 output is the
collection of \emph{MLE evaluation claims}
\[
  a_g \;=\; \mle{\psia(w_g)}(r^*) \;\in\; \GFt,
  \qquad g = 1,\ldots,K,
\]
at a common sumcheck endpoint $r^* \in \GFt^{\nu}$. The verifier
holds $(r^*, \alpha, a_1, \ldots, a_K)$ and the commitment root; the
prover holds the full witness and the codeword matrices from commit.
The goal of Step~7 is to certify each $a_g$ relative to the
commitment.

\paragraph{The eq-tensor.} For analysis we always assume the matrix
shape above, so that the eq-MLE factorises as a tensor product:
\[
  \eq_b(r^*) \;=\; q_0[b_0]\cdot q_1[b_1],
\]
where $q_0 \in \GFt^{n_{\mathrm{rows}}}$, $q_1 \in
\GFt^{n_{\mathrm{cols}}}$ are the two halves of the tensor
decomposition. They are computed identically by prover and verifier
from $r^*$ via the standard \texttt{build\_eq\_x\_r} routine. The
matrix form of the MLE evaluation claim is
\[
  a_g \;=\; q_0^{\top}\, \psia(M_{w_g})\, q_1 \;\in\; \GFt,
\]
where $\psia(M_{w_g})$ is the matrix with $\psia(M_{w_g})_{b_0, b_1}
= \psia(M_{w_g, (b_0, b_1)})$.

\paragraph{Linear code.} The commit phase encodes each witness column
row by row through an $\F_2$-linear code
\[
  \encode : \cellring^{n_{\mathrm{cols}}} \;\longrightarrow\;
            \cellring^{n_{\mathrm{cw}}}, \qquad n_{\mathrm{cw}} = 4\cdot n_{\mathrm{cols}},
\]
deployed in Zinc+ as the rate-$1/4$ \texttt{RaaF2Code} (Repeat /
permute / XOR-accumulate / permute / XOR-accumulate). Two facts about
this code matter for what follows:

\begin{lemma}[Linearity of the deployed code]
\label{lem:f2x-linearity}
The encoder $\encode : \cellring^{n_{\mathrm{cols}}} \to
\cellring^{n_{\mathrm{cw}}}$ is $\F_2$-linear in the cell
coordinate. Consequently it extends naturally to a
$\polyring$-linear map
$\encode : \polyring^{n_{\mathrm{cols}}} \to
\polyring^{n_{\mathrm{cw}}}$ acting cell-wise on any wider cell
ring $\polyringat{W} \supset \cellring$, and the diagram
\[
  \begin{array}{ccc}
    \polyringat{W}^{n_{\mathrm{cols}}} & \xrightarrow{\;\encode\;} & \polyringat{W}^{n_{\mathrm{cw}}} \\
    \downarrow{\scriptstyle c \cdot} & & \downarrow{\scriptstyle c \cdot} \\
    \polyringat{W}^{n_{\mathrm{cols}}} & \xrightarrow{\;\encode\;} & \polyringat{W}^{n_{\mathrm{cw}}}
  \end{array}
\]
commutes for every scalar $c \in \polyring$.
\end{lemma}

\begin{proof}
The kernel of \texttt{RaaF2Code} consists of four steps: integer
repetition (cell duplication into $\mathrm{REP}=4$ copies of the input
row), index permutation (no per-cell arithmetic), and pre-fix XOR
accumulation along the codeword. Repetition and permutation are
$\F_2$-linear in the cell coordinate trivially; XOR-accumulation is
$\F_2$-addition along the row, which is also $\F_2$-linear in the
cell coordinate. Composition of $\F_2$-linear maps is $\F_2$-linear.
Extending scalars from $\F_2$ to $\polyring$ preserves linearity
because $\polyring$ is a flat $\F_2$-module: scalar multiplication by
$c \in \polyring$ acts cell-wise, commutes with the XOR (which is
$\F_2$-addition on coefficients), and commutes with permutation
(which is positional). The commuting diagram follows.
\end{proof}

In particular \emph{the commit-time codeword matrix
$M_{\mathrm{cw}, g}$, originally encoded under
$\encode : \cellring^{n_{\mathrm{cols}}} \to
\cellring^{n_{\mathrm{cw}}}$, is also the codeword of the same row
under any wider $\encode : \polyringat{W}^{n_{\mathrm{cols}}} \to
\polyringat{W}^{n_{\mathrm{cw}}}$.} The commitment is not re-issued
for the Step~7 opening.
